\section{Síntesis y ampliación}

Esta clase presenta un marco de tres niveles para pasar de la precisión en benchmarks a la evaluación de agentes:

\begin{enumerate}
    \item \textbf{Capability span}: METR utiliza el tiempo de finalización humano para estimar cuánto tiempo puede sostener una tarea un modelo con una tasa de éxito dada;
    \item \textbf{Economic usefulness}: GDPval permite que expertos de la industria comparen directamente los productos de trabajo de la IA con los de profesionales, considerando además los costos de reintento, corrección y revisión;
    \item \textbf{Epistemic quality}: DeepScholar-Bench examina la cobertura de fuentes, la síntesis de conocimiento y la verificabilidad de las citas.
\end{enumerate}

La conclusión más importante no es que «los agentes ya pueden hacer tareas de una hora» ni que «se acercan a la tasa de victorias de los expertos», sino que el crecimiento de las capacidades muestra una fuerte condicionalidad:

\begin{itemize}
    \item las tareas largas con una tasa de éxito del 50\% ven su duración reducida de manera significativa al exigir una confiabilidad del 80\%;
    \item el desempeño en benchmarks con un prompt completo no representa la adquisición de contexto en un entorno real;
    \item la calidad visual, el seguimiento de instrucciones, la exactitud factual y la cobertura de citas pueden estar dominados por modelos distintos entre sí;
    \item una salida «aceptable pero no perfecta» puede generar valor en escenarios donde el costo de corrección es bajo, pero no es adecuada para despliegues de alto riesgo sin supervisión humana;
    \item el progreso en horizontes largos debe ir acompañado de una mejor gestión de estado, estrategias de detención, recuperación de errores y verifiers.
\end{itemize}

En definitiva, la evaluación no es una sola tabla de clasificación general, sino un sistema de diagnóstico para el entrenamiento, el enrutamiento y el control de riesgos. La siguiente clase regresará al tema central del curso: cómo superar los cuellos de botella de datos, verificación y eficiencia en la autosuperación, explorando el fine-tuning multiagente, la meta-verificación, Absolute Zero y la inteligencia por vatio.

\end{document}
